%\makeglossaries
% \makenoidxglossaries

% ===== Danh mục từ viết tắt =====
\newglossaryentry{adrmse}{
    type=\acronymtype, name={ADR-MSE},
    description={Approx Discounted Rank MSE (hàm mất mát chưng cất dựa trên thứ hạng)},
    first={Approx Discounted Rank MSE (ADR-MSE)}
}
\newglossaryentry{api}{
    type=\acronymtype, name={API},
    description={Giao diện lập trình ứng dụng (Application Programming Interface)},
    first={Giao diện lập trình ứng dụng (Application Programming Interface -- API)}
}
\newglossaryentry{bert}{
    type=\acronymtype, name={BERT},
    description={Bidirectional Encoder Representations from Transformers},
    first={Bidirectional Encoder Representations from Transformers (BERT)}
}
\newglossaryentry{bm25}{
    type=\acronymtype, name={BM25},
    description={Thuật toán xếp hạng văn bản theo từ khóa (Best Matching 25)},
    first={Best Matching 25 (BM25)}
}
\newglossaryentry{cntt}{
    type=\acronymtype, name={CNTT},
    description={Công nghệ thông tin},
    first={Công nghệ thông tin (CNTT)}
}
\newglossaryentry{cot}{
    type=\acronymtype, name={CoT},
    description={Chuỗi suy luận từng bước (Chain-of-Thought)},
    first={Chain-of-Thought (CoT)}
}
\newglossaryentry{cpu}{
    type=\acronymtype, name={CPU},
    description={Bộ xử lý trung tâm (Central Processing Unit)},
    first={Central Processing Unit (CPU)}
}
\newglossaryentry{datn}{
    type=\acronymtype, name={ĐATN},
    description={Đồ án tốt nghiệp},
    first={Đồ án tốt nghiệp (ĐATN)}
}
\newglossaryentry{gist}{
    type=\acronymtype, name={GIST},
    description={Guided In-sample Selection of Training Negatives (chọn lọc negative có hướng dẫn)},
    first={Guided In-sample Selection of Training Negatives (GIST)}
}
\newglossaryentry{gpu}{
    type=\acronymtype, name={GPU},
    description={Bộ xử lý đồ họa (Graphics Processing Unit)},
    first={Graphics Processing Unit (GPU)}
}
\newglossaryentry{kd}{
    type=\acronymtype, name={KD},
    description={Chưng cất tri thức (Knowledge Distillation)},
    first={Knowledge Distillation (KD)}
}
\newglossaryentry{llm}{
    type=\acronymtype, name={LLM},
    description={Mô hình ngôn ngữ lớn (Large Language Model)},
    first={Large Language Model (LLM)}
}
\newglossaryentry{mcq}{
    type=\acronymtype, name={MCQ},
    description={Câu hỏi trắc nghiệm (Multiple Choice Question)},
    first={Multiple Choice Question (MCQ)}
}
\newglossaryentry{mmarco}{
    type=\acronymtype, name={mMARCO},
    description={Bộ dữ liệu MS MARCO đa ngôn ngữ (multilingual MS MARCO)},
    first={multilingual MS MARCO (mMARCO)}
}
\newglossaryentry{mnr}{
    type=\acronymtype, name={MNR},
    description={Hàm mất mát xếp hạng đa negative (Multiple Negatives Ranking Loss)},
    first={Multiple Negatives Ranking Loss (MNR)}
}
\newglossaryentry{mrl}{
    type=\acronymtype, name={MRL},
    description={Học biểu diễn lồng nhau (Matryoshka Representation Learning)},
    first={Matryoshka Representation Learning (MRL)}
}
\newglossaryentry{mrr}{
    type=\acronymtype, name={MRR},
    description={Nghịch đảo thứ hạng trung bình (Mean Reciprocal Rank)},
    first={Mean Reciprocal Rank (MRR)}
}
\newglossaryentry{ndcg}{
    type=\acronymtype, name={NDCG},
    description={Độ lợi tích lũy chiết khấu chuẩn hóa (Normalized Discounted Cumulative Gain)},
    first={Normalized Discounted Cumulative Gain (NDCG)}
}
\newglossaryentry{pdf}{
    type=\acronymtype, name={PDF},
    description={Định dạng tài liệu di động (Portable Document Format)},
    first={Portable Document Format (PDF)}
}
\newglossaryentry{rag}{
    type=\acronymtype, name={RAG},
    description={Sinh có tăng cường truy hồi (Retrieval-Augmented Generation)},
    first={Retrieval-Augmented Generation (RAG)}
}
\newglossaryentry{rrf}{
    type=\acronymtype, name={RRF},
    description={Hợp nhất theo thứ hạng nghịch đảo (Reciprocal Rank Fusion)},
    first={Reciprocal Rank Fusion (RRF)}
}
\newglossaryentry{sv}{
    type=\acronymtype, name={SV},
    description={Sinh viên},
    first={Sinh viên (SV)}
}
\newglossaryentry{vram}{
    type=\acronymtype, name={VRAM},
    description={Bộ nhớ truy cập ngẫu nhiên của GPU (Video Random Access Memory)},
    first={Video Random Access Memory (VRAM)}
}